% SPDX-License-Identifier: Apache-2.0
\section{What Rust Cannot Do}
\label{sec:e-walls}

One wall is left, and it is the one a separate parser was buying.

\subsection{Control flow on a signal}

A signal is a node in a graph. It has no value while the Rust program
runs, so it cannot be a condition.

\begin{lstlisting}[caption={The wall. A signal has no value while the program runs, so it cannot be a condition.},label={lst:wall}]
if cond { a } else { b }   // cond must be a bool.
                           // A signal is not a bool.
\end{lstlisting}

\code{if} and \code{match} take a \code{bool} and a pattern, and neither
can be overloaded. Control flow on a signal needs another spelling.

\subsubsection{The name is not \code{if\_!}}

The obvious macro name is \code{if\_!}, and it should be rejected for a
better reason than the trailing underscore. \code{if} branches: one arm
runs and the other does not. This does not branch. Both arms exist in the
hardware at once and a multiplexer picks between them, so area is paid for
both and a side effect in the arm not chosen still has to be suppressed. A
reader who brings Rust's \code{if} semantics to a construct named after it
is misled about the one thing that matters most.

\code{when!} is what Chisel~\cite{bachrach2012} and SpinalHDL call it, so
the engineers who will read it already know what it means, and it claims
nothing it does not do.

\subsubsection{Two forms, and where the limit is}

An expression needs no macro at all. \code{mux} is the operation, and a
function says so.

\begin{lstlisting}[caption={The expression form is a function. There is nothing to predicate.},label={lst:mux}]
let chosen = mux(reset, Sig(0), self.count.get());
\end{lstlisting}

A statement form has to predicate whatever is written inside it, and that
is what real logic needs.

\begin{lstlisting}[caption={The statement form. It reads as a conditional and lowers to predicated writes.},label={lst:when}]
when!(enable => {
    self.count => Sig(self.count.get().0 + 1);
    self.flag  => Sig(1)
} else {
    self.count => Sig(0);
    self.flag  => Sig(0)
});
\end{lstlisting}

Probe~10 compiles both. The limit is in the second: a
\code{macro\_rules!} macro matches a grammar, so this one predicates a
list of register writes and nothing else. It cannot reach inside an
arbitrary block and rewrite the statements it finds there, because
\code{macro\_rules!} cannot see into a block it did not parse itself.

Widening that grammar case by case ends up specifying a small language
inside the macro, which is the thing the embedding set out to avoid. A
proc macro reads the item's real syntax tree and can reinterpret a plain
\code{if}, which is why Section~\ref{sec:e-elaboration} recommends one.
This is the strongest single argument for that recommendation, and it
arrived from writing the macro rather than from reasoning about it.

\subsection{Bit slicing}

\code{w[0..8]} cannot work, because \code{Index} returns a reference to
something that already exists and a slice of a signal is a new node. It
becomes \code{w.slice::<0, 8>()}.

The bounds do not all have to be literals, and the rule is not the one a
first reading suggests.

\begin{table}[H]
\caption{What a slice bound may be. The width is the type, so it is the
one that cannot vary.}
\label{tab:slice}
\centering
\footnotesize
\begin{tabular}{@{}lll@{}}
\toprule
Form & Offset & Width \\
\midrule
\code{w.slice::<0, 8>()}        & literal          & literal \\
\code{w.slice::<LO, LEN>()}     & const parameter  & const parameter \\
\code{w.slice\_at::<8>(i)}      & \textbf{run time} & static \\
\code{w.slice\_var(i, n)}       & run time         & \textbf{rejected} \\
\bottomrule
\end{tabular}
\end{table}

A generic function may slice without knowing the numbers, because a const
generic parameter is a compile-time value and not a literal. An offset may
be a run-time value, which synthesises to a barrel shifter feeding a
fixed-width slice.

A width may not. The width is the type, and a type cannot depend on a
value that exists only at run time: probe~11b gives
\code{error[E0435]: attempt to use a non-constant value in a constant}.
That is not a Rust limitation to work around. A signal whose bit width is
decided at run time is not a thing hardware has, so the compiler is
refusing to express something that could not be built.

\section{What the Embedding Costs}
\label{sec:e-costs}

\subsection{Structural typing on data changes shape}

The merged specification resolved one of its four conflicts by making
transaction compatibility structural: two transactions with matching fields
are compatible whatever they are called~\cite{unification}. Rust is
nominally typed.

A derive can recover it by generating an associated layout type, so that
compatibility becomes a bound rather than a name comparison. It recovers it
positionally, though, so the rule becomes ``the same field types in the
same order'' rather than ``the same fields''. That is a real change to a
decision already taken, and unlike everything else in this paper it has not
been compiled.

\subsection{A thesis is inverted}

The project's second thesis~\cite{theses} names making the language a valid
program in an existing language as one direction, and then takes the other,
on the grounds that a full toolchain has become affordable to build. This
paper takes the direction the thesis declined.

The thesis is the rationale section of the merged specification, so it
should be rewritten rather than quietly contradicted. Its argument survives
the change of direction, and only its conclusion moves: a toolchain is
still being built, but it buys a front end over Rust's own syntax tree
rather than a lexer and a parser, and the type system, the module system,
the generics and the tooling arrive already written.

\subsection{Smaller frictions}

\code{concat!} is a \code{std} macro, so the bit concatenation macro needs
another name. \code{async fn} in a public trait warns that auto trait
bounds cannot be stated, which is harmless while elaboration is single
threaded and would matter if it ever were not.
